\documentclass[12pt,a4paper]{article}

% ============================================================
% PACOTES
% ============================================================
\usepackage[utf8]{inputenc}
\usepackage[T1]{fontenc}
\usepackage[brazil]{babel}
\usepackage{tikz}
\usepackage{pgfplots}
\pgfplotsset{compat=1.18}

\usepackage{amsmath,amssymb,amsfonts,mathtools}
\usepackage{bm}
\usepackage{booktabs}
\usepackage{array}
\usepackage{geometry}
\usepackage{graphicx}
\usepackage{float}
\usepackage{xcolor}
\usepackage{hyperref}
\usepackage{microtype}

\geometry{margin=2.5cm}

\hypersetup{
    colorlinks=true,
    linkcolor=blue,
    citecolor=blue,
    urlcolor=blue
}

% ============================================================
% COMANDOS
% ============================================================
\newcommand{\e}{\mathrm{e}}
\newcommand{\dd}{\mathrm{d}}
\newcommand{\avg}[1]{\left\langle #1\right\rangle}
\newcommand{\Tr}{\operatorname{Tr}}
\newcommand{\sech}{\operatorname{sech}}
\newcommand{\Var}{\operatorname{Var}}
\newcommand{\Prob}{\mathbb{P}}
\newcommand{\sgn}{\operatorname{sgn}}
\newcommand{\erf}{\operatorname{erf}}
\newcommand{\D}{\mathcal{D}}

% ============================================================
% DOCUMENTO
% ============================================================
\begin{document}
%
\begin{center}
    {\Large\bfseries Sétima Série de Exercícios}\\[0.6cm]
    {\large Tópicos de Mecânica Estatística -- Transições de Fases -- 2026}\\[0.2cm]
    {\large Campos aleatórios e vidros de spin}\\[0.2cm]
    {\large Data de Entrega -- 13/11/2026}\\[0.2cm]
    {\Large \textbf{Autor:} Thiago Siqueira Domingues (Nº USP: 8942194)}
\end{center}
% ============================================================
% ============================================================
\section*{Problema 1 -- Modelo de Ising com campos aleatórios}
% ============================================================

Considere o modelo de Ising totalmente conectado com Hamiltoniano
%
\begin{equation}
    \mathcal{H}
    =
    -\frac{J}{2N}
    \left(\sum_{i}\sigma_i\right)^2
    -\sum_{i}H_i\sigma_i,
    \qquad \sigma_i=\pm1,\quad J>0,
    \label{eq:rfim_H}
\end{equation}
%
onde os campos $H_i$ são variáveis aleatórias \emph{quenched},
independentes e identicamente distribuídas.

%------------------------------------------------------------
\subsection*{(i) Distribuição bimodal}
%------------------------------------------------------------

Consideramos $p(H_i)=\tfrac{1}{2}[\delta(H_i-H)+\delta(H_i+H)]$.
Pela transformação de Hubbard--Stratonovich, a função de partição
no limite $N\to\infty$ é dominada pelo mínimo da energia livre
efetiva por sítio~\cite{Aharony1978a}:
%
\begin{equation}
    f(m)
    =
    \frac{Jm^2}{2}
    -\frac{1}{2\beta}\ln\bigl[2\cosh\beta(Jm+H)\bigr]
    -\frac{1}{2\beta}\ln\bigl[2\cosh\beta(Jm-H)\bigr].
    \label{eq:free_energy_bimodal}
\end{equation}
%
A condição $\partial f/\partial m=0$ fornece a equação de autoconsistência
%
\begin{equation}
    \boxed{
    m
    =
    \frac{1}{2}\bigl[
        \tanh\beta(Jm+H)
        +
        \tanh\beta(Jm-H)
    \bigr].
    }
    \label{eq:self_consistency_bimodal}
\end{equation}
%
Por simetria, $m=0$ é sempre solução; a simetria $m\to-m$ garante
que a magnetização espontânea emerge continuamente ou descontinuamente.

\paragraph{Linha crítica contínua.}
Expandindo $f(m)$ em torno de $m=0$, a instabilidade do estado
paramagnético ocorre quando $\partial^2 f/\partial m^2|_{m=0}=0$, ou seja,
%
\begin{equation}
    \beta J\sech^2(\beta H)=1.
    \label{eq:continuous_bimodal}
\end{equation}
%
Introduzindo $x\equiv H/k_BT$, a linha crítica admite a parametrização exata
%
\begin{equation}
    t\equiv\frac{k_BT}{J}=\sech^2 x,
    \qquad
    h\equiv\frac{H}{J}=x\sech^2 x.
    \label{eq:critical_line}
\end{equation}
%
Para $H=0$ ($x=0$) recuperamos a temperatura crítica do modelo de Ising puro:
$k_BT_c=J$.

\paragraph{Ponto tricrítico.}
Para investigar a natureza da transição, expandimos
\eqref{eq:self_consistency_bimodal} em potências de $m$
com $y=\beta Jm$ e $x=\beta H$ fixo:
%
\begin{equation}
    0
    =
    \underbrace{\bigl(\beta J\sech^2 x-1\bigr)}_{A}\,m
    +
    \underbrace{\frac{(\beta J)^3}{3}\sech^2 x\,(3\tanh^2 x-1)}_{-B}\,m^3
    +O(m^5).
    \label{eq:landau_equation}
\end{equation}
%
A transição é de segunda ordem quando $A=0$ com $B>0$, e de primeira ordem
quando $B$ muda de sinal. O ponto tricrítico (TCP) ocorre em $A=B=0$
simultaneamente, impondo
%
\begin{equation}
    \tanh^2 x_{\mathrm{tc}}=\frac{1}{3}
    \;\Longrightarrow\;
    \sech^2 x_{\mathrm{tc}}=\frac{2}{3}.
    \label{eq:tcp_condition}
\end{equation}
%
Substituindo na condição crítica \eqref{eq:continuous_bimodal}:
%
\begin{equation}
    \boxed{
    \frac{k_BT_{\mathrm{tc}}}{J}=\frac{2}{3},
    \qquad
    \frac{H_{\mathrm{tc}}}{J}
    =\frac{2}{3}\operatorname{artanh}\!\left(\frac{1}{\sqrt{3}}\right)
    \simeq 0.439.
    }
    \label{eq:tcp}
\end{equation}

\paragraph{Linha de primeira ordem e limite $T\to0$.}
Abaixo do TCP, a fronteira de equilíbrio satisfaz simultaneamente
a equação de autoconsistência e a condição de coexistência de fases
$f(m_\star)=f(0)$, que se escreve explicitamente
%
\begin{equation}
    \frac{Jm_\star^2}{2}
    =
    \frac{1}{2\beta}
    \ln\!\left[
        \frac{
        \cosh\beta(Jm_\star+H)\,\cosh\beta(Jm_\star-H)
        }{
        \cosh^2(\beta H)
        }
    \right].
    \label{eq:coexistence}
\end{equation}
%
Este sistema determina numericamente a linha de primeira ordem.
A $T=0$, a comparação direta de energias --- estado ferromagnético
($e_F=-J/2$) versus estado com cada spin alinhado ao seu campo
($e_R=-H$) --- dá a fronteira
%
\begin{equation}
    \boxed{H_c(T=0)=\frac{J}{2}.}
\end{equation}
%
O diagrama de fases completo é mostrado na Fig.~\ref{fig:rfim_phase_diagram}.

\begin{figure}[H]
\centering
\begin{tikzpicture}[scale=5.5, >=stealth]

  % eixos
  \draw[->] (-0.02,0) -- (1.08,0) node[right, font=\small] {$H/J$};
  \draw[->] (0,-0.02) -- (0,1.12) node[above, font=\small] {$k_BT/J$};
  \node[below left, font=\small] at (0,0) {$0$};

  % marcações
  \draw (0.439, 0.012) -- (0.439,-0.012);
  \node[below, font=\tiny] at (0.439,0) {$H_{\rm tc}/J$};
  \draw (0.5, 0.012) -- (0.5,-0.012);
  \node[below, font=\tiny] at (0.5,0) {$\tfrac{1}{2}$};
  \draw (1.0, 0.012) -- (1.0,-0.012);
  \node[below, font=\tiny] at (1.0,0) {$1$};
  \draw (0.010,0.667) -- (-0.010,0.667);
  \node[left, font=\tiny] at (0,0.667) {$\tfrac{2}{3}$};
  \draw (0.010,1.0) -- (-0.010,1.0);
  \node[left, font=\tiny] at (0,1.0) {$1$};

  % linha de segunda ordem: t = sech^2(x), h = x sech^2(x)
  \draw[thick, samples=200, domain=0:0.6585,
        /pgf/fpu, /pgf/fpu/output format=fixed]
    plot ({(\x)/cosh(\x)/cosh(\x)}, {1/cosh(\x)/cosh(\x)});

  % linha de primeira ordem (amostrada numericamente)
  \draw[thick, dashed]
    plot[smooth] coordinates {
      (0.439, 0.667)
      (0.455, 0.570)
      (0.470, 0.460)
      (0.482, 0.340)
      (0.492, 0.210)
      (0.498, 0.100)
      (0.500, 0.000)
    };

  % pontos especiais
  \filldraw (0,1) circle (0.4pt);
  \node[right, font=\footnotesize] at (0.008,1.05) {$T_c$};
  \filldraw (0.439,0.667) circle (0.4pt);
  \node[above right, font=\footnotesize] at (0.439,0.667) {TCP};
  \filldraw (0.500,0) circle (0.4pt);

  % rótulos de fase
  \node[font=\normalsize\bfseries] at (0.12, 0.76) {FM};
  \node[font=\normalsize\bfseries] at (0.72, 0.76) {PM};

  % anotações
  \node[font=\scriptsize] at (0.19, 0.50) {\rotatebox{-52}{2ª ordem}};
  \node[font=\scriptsize] at (0.56, 0.32) {\rotatebox{-80}{1ª ordem}};

\end{tikzpicture}
\caption{Diagrama de fases do RFIM de campo médio com distribuição
bimodal, reproduzindo a Fig.~1 de Aharony~\cite{Aharony1978a}.
FM: fase ferromagnética; PM: fase paramagnética.
A linha sólida é a fronteira de segunda ordem
[$k_BT/J = \sech^2(H/k_BT)$, Eq.~\eqref{eq:critical_line}],
que termina no ponto tricrítico TCP $\approx(0.439,\,0.667)$.
A linha tracejada é de primeira ordem e alcança $T=0$ em $H/J=1/2$.}
\label{fig:rfim_phase_diagram}
\end{figure}
%
%------------------------------------------------------------
\subsection*{(ii) Distribuição gaussiana}
%------------------------------------------------------------

Consideramos agora $p(H_i) = (2\pi H^2)^{-1/2}\exp(-H_i^2/2H^2)$,
com $H$ o desvio-padrão. A equação de autoconsistência torna-se
%
\begin{equation}
    m
    =
    \int_{-\infty}^{\infty}\dd H_i\;p(H_i)\,
    \tanh\bigl[\beta(Jm+H_i)\bigr].
    \label{eq:self_consistency_gaussian}
\end{equation}
%
Expandindo para $m$ pequeno e usando $z=H_i/H$, a condição de
instabilidade do estado paramagnético é~\cite{SchneiderPytte1977}
%
\begin{equation}
    \boxed{
    1
    =
    \frac{1}{t\sqrt{2\pi}}
    \int_{-\infty}^{\infty}\dd z\;
    \e^{-z^2/2}
    \sech^2\!\left(\frac{d}{t}z\right),
    }
    \label{eq:critical_gaussian}
\end{equation}
%
onde $t=k_BT/J$ e $d=H/J$.
Esta equação define implicitamente a linha crítica — sempre contínua,
pois o coeficiente $B$ da expansão de Landau satisfaz $B>0$ em todo o
diagrama para a distribuição gaussiana~\cite{Aharony1978a}.

\paragraph{Limite $T\to0$.}
Quando $\beta\to\infty$, $\tanh[\beta(Jm+H_i)]\to\sgn(Jm+H_i)$,
e a equação de estado \eqref{eq:self_consistency_gaussian} reduz-se a
%
\begin{equation}
    m = \erf\!\left(\frac{Jm}{\sqrt{2}\,H}\right),
    \label{eq:gaussian_T0}
\end{equation}
%
onde $\erf(x)=\frac{2}{\sqrt{\pi}}\int_0^x\e^{-u^2}\dd u$.
Expandindo para $m$ pequeno: $\erf(x)\approx\frac{2}{\sqrt{\pi}}x$,
a condição de instabilidade de $m=0$ a $T=0$ fornece
%
\begin{equation}
    \boxed{
    \frac{H_c(0)}{J} = \sqrt{\frac{2}{\pi}} \simeq 0.798.
    }
    \label{eq:gaussian_Hc0}
\end{equation}

\paragraph{Comparação com a distribuição bimodal.}
A ausência de ponto tricrítico na distribuição gaussiana tem origem
no sinal do coeficiente de Landau $B$. Aharony~\cite{Aharony1978a}
mostrou que $B \propto -p''(0)$: para a gaussiana $p''(0)<0$
implica $B>0$ em todo o diagrama, garantindo transição contínua.
Para a distribuição bimodal, $p''(0)>0$ para campos suficientemente
intensos, tornando $B<0$ e gerando a transição de primeira ordem
com ponto tricrítico. A diferença é portanto ditada pela curvatura
da distribuição em $H=0$. O diagrama de fases é mostrado na Fig.~\ref{fig:rfim_gaussian_phase_diagram}.
%
\begin{figure}[H]
\centering
\begin{tikzpicture}[scale=5.5, >=stealth]

  % eixos
  \draw[->] (-0.02,0) -- (1.00,0)
    node[right, font=\small] {$H/J$};
  \draw[->] (0,-0.02) -- (0,1.12)
    node[above, font=\small] {$k_BT/J$};
  \node[below left, font=\small] at (0,0) {$0$};

  % marcações
  \draw (0.798, 0.012) -- (0.798,-0.012);
  \node[below, font=\small] at (0.798,0) {$\sqrt{2/\pi}$};
  \draw (0.010,1.0) -- (-0.010,1.0);
  \node[left, font=\small] at (0,1.0) {$1$};

  % linha crítica contínua (pontos numéricos)
  \draw[thick]
    plot[smooth, tension=0.6] coordinates {
      (0.000, 1.000)
      (0.072, 0.950)
      (0.145, 0.900)
      (0.213, 0.850)
      (0.277, 0.800)
      (0.336, 0.750)
      (0.391, 0.700)
      (0.440, 0.650)
      (0.484, 0.600)
      (0.523, 0.550)
      (0.557, 0.500)
      (0.587, 0.450)
      (0.613, 0.400)
      (0.634, 0.350)
      (0.652, 0.300)
      (0.667, 0.250)
      (0.678, 0.200)
      (0.688, 0.150)
      (0.694, 0.100)
      (0.697, 0.050)
      (0.798, 0.000)
    };

  % pontos especiais
  \filldraw (0,1) circle (0.4pt);
  \node[right, font=\footnotesize] at (0.008,1.0) {$T_c$};
  \filldraw (0.798,0) circle (0.4pt);

  % rótulos de fase
  \node[font=\normalsize\bfseries] at (0.18, 0.60) {FM};
  \node[font=\normalsize\bfseries] at (0.72, 0.70) {PM};

  % anotação
  \node[font=\scriptsize] at (0.42, 0.44)
    {\rotatebox{-60}{2ª ordem}};

\end{tikzpicture}
\caption{Diagrama de fases do RFIM de campo médio com distribuição
gaussiana, seguindo Schneider e Pytte~\cite{SchneiderPytte1977}. A convenção de Schneider e Pytte~\cite{SchneiderPytte1977} usa os eixos invertidos ($2J/\sigma$ vs.\ $k_BT/\sigma$),
o que transforma a curva crescente aqui em decrescente, mas
representa o mesmo diagrama físico. FM: fase ferromagnética; PM: fase paramagnética. A fronteira é sempre de segunda ordem [Eq.~\eqref{eq:critical_gaussian}],
parte de $(0,1)$ e termina em $H_c(0)/J=\sqrt{2/\pi}\approx0.798$
a $T=0$ [Eq.~\eqref{eq:gaussian_Hc0}].
Não há ponto tricrítico.}
\label{fig:rfim_gaussian_phase_diagram}
\end{figure}
%
%------------------------------------------------------------
\subsection*{(iii) Método das réplicas}
%------------------------------------------------------------

Para desordem \emph{quenched}, a grandeza termodinâmica relevante é
a energia livre média $\overline{F} = -\beta^{-1}\overline{\ln Z}$.
O cálculo direto de $\overline{\ln Z}$ é difícil; o método das
réplicas contorna isso pela identidade
%
\begin{equation}
    \overline{\ln Z}
    =
    \lim_{n\to 0}\frac{\overline{Z^n}-1}{n}.
    \label{eq:replica_identity}
\end{equation}
%
Introduzimos $n$ cópias do sistema com a mesma realização de $\{H_i\}$.
A função de partição replicada é
%
\begin{equation}
    Z^n
    =
    \sum_{\{\sigma_i^a\}}
    \exp\!\left[
        \frac{\beta J}{2N}
        \sum_{a=1}^{n}
        \!\left(\sum_{i}\sigma_i^a\right)^{\!2}
        +
        \beta\sum_{i}H_i\sum_{a}\sigma_i^a
    \right].
    \label{eq:Zn}
\end{equation}
%
Realizando a média sobre os campos, que são independentes entre sítios,
%
\begin{equation}
    \overline{Z^n}
    =
    \sum_{\{\sigma_i^a\}}
    \exp\!\left[
        \frac{\beta J}{2N}
        \sum_{a}\!\left(\sum_i\sigma_i^a\right)^{\!2}
    \right]
    \prod_{i}
    \overline{\exp\!\left(\beta H_i\sum_a\sigma_i^a\right)}.
    \label{eq:Zn_avg}
\end{equation}

\paragraph{Distribuição bimodal.}
Para $p(H_i)=\frac{1}{2}[\delta(H_i-H)+\delta(H_i+H)]$:
%
\begin{equation}
    \overline{\exp\!\left(\beta H_i\sum_a\sigma_i^a\right)}
    =
    \cosh\!\left(\beta H\sum_a\sigma_i^a\right).
    \label{eq:avg_bimodal}
\end{equation}
%
Desacoplando a interação ferromagnética por Hubbard--Stratonovich
e impondo a condição de ponto de sela na ação efetiva com o ansatz
simétrico $m_a = m$, recupera-se a equação de autoconsistência
\eqref{eq:self_consistency_bimodal}.

\paragraph{Distribuição gaussiana.}
Para $p(H_i)=\mathcal{N}(0,H^2)$, a média é o gerador de momentos:
%
\begin{equation}
    \overline{\exp\!\left(\beta H_i\sum_a\sigma_i^a\right)}
    =
    \exp\!\left[\frac{\beta^2 H^2}{2}\!\left(\sum_a\sigma_i^a\right)^{\!2}\right].
    \label{eq:avg_gaussian}
\end{equation}
%
O termo quadrático em $\sum_a\sigma_i^a$ acopla as réplicas entre si
e induz um parâmetro de ordem de vidro de spins
$q_{\rm EA} = \overline{\langle\sigma_i\rangle^2}$.
Contudo, Schneider e Pytte~\cite{SchneiderPytte1977} mostraram que,
para o RFIM com campo médio infinito, a solução de réplicas é
\emph{exata}: não há quebra de simetria entre réplicas, e o parâmetro
$q_{\rm EA}$ descreve a fase de spins independentes ($m=0$, $q_{\rm EA}\neq0$),
não uma fase de vidro de spins com correlações.
No ansatz simétrico $m_a=m$, a condição estacionária reproduz
a equação de autoconsistência \eqref{eq:self_consistency_gaussian}.

\paragraph{Conclusão.}
O método das réplicas recupera, para ambas as distribuições, as equações
de campo médio dos itens (i) e (ii). A diferença essencial é que, para
campos gaussianos, a média \eqref{eq:avg_gaussian} gera um acoplamento
inter-réplica $(\sum_a\sigma_i^a)^2$, o que pode, em princípio,
sinalizar uma fase de vidro de spins. Para o RFIM, porém, esse acoplamento
não implica correlações entre spins — a fase desordenada é de spins
\emph{independentes} alinhados pelos campos locais.
%
\section*{Problema 2 -- Modelo de vidros de spin de van Hemmen}
% ============================================================
Vou seguir a construção de van Hemmen~\cite{vanHemmen1982}, completando
os passos algébricos que o artigo deixa em aberto.

\subsection*{Hamiltoniano e variáveis macroscópicas}
O Hamiltoniano de $N$ spins de Ising, sem campo externo, é
%
\begin{equation}
    \mathcal{H}_N
    =
    -\frac{J_0}{N}\sum_{i<j}\sigma_i\sigma_j
    -\sum_{i<j}J_{ij}\sigma_i\sigma_j,
    \qquad
    J_{ij}=\frac{J}{N}\left(\xi_i\eta_j+\xi_j\eta_i\right),
    \label{eq:vh_H}
\end{equation}
%
com $J>0$, $J_0\ge0$. As variáveis $\xi_i,\eta_i$ são independentes, de
média nula e variância unitária. O ponto-chave é que $\mathcal{H}_N$
depende dos spins apenas através de três somas. Definimos
%
\begin{equation}
    m_N=\frac1N\sum_i\sigma_i,\qquad
    q_{1N}=\frac1N\sum_i\xi_i\sigma_i,\qquad
    q_{2N}=\frac1N\sum_i\eta_i\sigma_i .
    \label{eq:vh_vars}
\end{equation}
%
(Aqui $m_N$, $q_{1N}$ e $q_{2N}$ são funções da configuração de spins;
as médias térmicas só aparecem depois, no limite $N\to\infty$.)
Como $\sum_{i<j}\sigma_i\sigma_j=\tfrac12(N^2m_N^2-N)$ e
$\sum_{i<j}J_{ij}\sigma_i\sigma_j=J N q_{1N}q_{2N}+O(1)$, resulta, com
$K_0=\beta J_0$ e $K=\beta J$,
%
\begin{equation}
    -\beta\mathcal{H}_N
    =
    N\left(\tfrac12K_0m_N^2+Kq_{1N}q_{2N}\right)
    \equiv N\,Q(\mathbf{m}_N),
    \qquad
    \mathbf{m}_N=(m_N,q_{1N},q_{2N}).
    \label{eq:vh_Q}
\end{equation}
%
Em vez de somar sobre $2^N$ configurações, passamos a somar sobre o
vetor $\mathbf{m}$ de três componentes; falta o ``jacobiano'' dessa
mudança de variáveis, que é a entropia associada a cada valor de $\mathbf{m}$.

\subsection*{Entropia por grandes desvios}
Seja $E_s(X)=2^{-N}\mathrm{Tr}(X)$ a média sobre configurações de spin
(spins independentes, $\pm1$ equiprováveis). A probabilidade de
$\mathbf{m}_N$ estar perto de $\mathbf{m}$ é, pela teoria de grandes
desvios (Cramér), $\exp[-N c^*(\mathbf{m})]\,d\mathbf{m}$, onde
$c^*$ é a transformada de Legendre
%
\begin{equation}
    c(\mathbf{t})=\lim_{N\to\infty}\frac1N\ln E_s\!\left(e^{\mathbf{t}\cdot N\mathbf{m}_N}\right)
    =\avg{\ln\cosh(t_0+t_1\xi+t_2\eta)}_{\xi,\eta},
    \qquad
    c^*(\mathbf{m})=\sup_{\mathbf{t}}\left[\mathbf{m}\cdot\mathbf{t}-c(\mathbf{t})\right].
    \label{eq:vh_c}
\end{equation}
%
A função $c$ é convexa e $c^*\ge0$, e $-c^*$ é a entropia média (a menos
de $\ln2$ da normalização). A energia livre por sítio é então
%
\begin{equation}
    -\beta f
    =\lim_{N\to\infty}\frac1N\ln\mathrm{Tr}\,e^{-\beta\mathcal{H}_N}
    =\ln2+\max_{\mathbf{m}}\left[Q(\mathbf{m})-c^*(\mathbf{m})\right].
    \label{eq:vh_f_max}
\end{equation}
%
\subsection*{Equações de autoconsistência}
Pela convexidade de $c^*$, o máximo em \eqref{eq:vh_f_max} ocorre em
$q_1=q_2\equiv q$. No máximo, $\nabla Q=\nabla c^*=\mathbf{t}$, ou seja
$t_0=K_0m$, $t_1=Kq_2$, $t_2=Kq_1$, e $\nabla c=\mathbf{m}$ fornece
%
\begin{equation}
    m=\avg{\tanh\!\big[K_0m+Kq(\xi+\eta)\big]}_{\xi,\eta},\qquad
    q=\avg{\tfrac12(\xi+\eta)\tanh\!\big[K_0m+Kq(\xi+\eta)\big]}_{\xi,\eta}.
    \label{eq:vh_sc}
\end{equation}
%
Para $\xi,\eta=\pm1$ com probabilidade $\tfrac12$ (a distribuição
bimodal usada no artigo para o diagrama de fases), a média tem quatro termos e
%
\begin{equation}
    \begin{aligned}
    m&=\tfrac14\left[\tanh(K_0m+2Kq)+2\tanh(K_0m)+\tanh(K_0m-2Kq)\right],\\
    q&=\tfrac14\left[\tanh(K_0m+2Kq)-\tanh(K_0m-2Kq)\right].
    \end{aligned}
    \label{eq:vh_sc4}
\end{equation}
%
A energia livre nos pontos estacionários é
$f=\tfrac12J_0m^2+Jq^2-\beta^{-1}\avg{\ln2\cosh[\beta(J_0m+J(\xi+\eta)q)]}$;
entre soluções distintas vale a de menor $f$.

\subsection*{Fases}
As soluções especiais de \eqref{eq:vh_sc4} são três (mais a mista):
paramagneto \textbf{P} ($m=q=0$); ferromagneto \textbf{F} ($m\neq0$,
$q=0$); vidro de spins \textbf{SG} ($m=0$, $q\neq0$); e fase mista
\textbf{M} ($m\neq0$, $q\neq0$).

\paragraph{Interpretação por cores.}
Como $\xi_i\eta_i=\pm1$, dividimos os sítios em \emph{azuis}
($\xi_i\eta_i=+1$) e \emph{vermelhos} ($\xi_i\eta_i=-1$). Se $i,j$ têm
cores diferentes, $\xi_i\eta_j+\xi_j\eta_i=\xi_i\eta_j(1+\xi_i\eta_i\xi_j\eta_j)\!\cdot\!\ldots=0$:
a interação aleatória só liga sítios de mesma cor, enquanto a
ferromagnética é cega à cor. Para $J_0=0$ os dois subsistemas desacoplam;
após a transformação de Mattis $\sigma_i\to\xi_i\sigma_i$ o azul vira
ferromagneto (ordena em $k_BT=J$) e o vermelho, antiferromagneto
totalmente frustrado, permanece desordenado até $T=0$. Com $J_0>0$ o
vermelho pode ordenar ferromagneticamente, ganhando energia pelo termo
$J_0$, o que dá a fase mista, com $K_0=\beta J_0\approx2$.

\subsection*{Diagrama de fases}
Em $(\alpha,t)=(J_0/J,\;k_BT/J)$ as fronteiras são:
%
\begin{itemize}
    \item \textbf{P--SG:} $t=1$ para $\alpha<1$ (de \eqref{eq:vh_sc4}, $q\simeq Kq$). Para $\alpha>1$, P--F em $t=\alpha$.
    \item \textbf{SG--M} (contínua): linearizando $m$ em torno da solução SG,
    %
    \begin{equation}
        t=\frac{\alpha}{2}\left[1+\mathrm{sech}^2\!\left(\frac{2q}{t}\right)\right],
        \qquad q=\tfrac12\tanh\!\left(\frac{2q}{t}\right),
        \label{eq:vh_SGM}
    \end{equation}
    %
    que vai de $t\simeq\alpha/2$ em $T\to0$ (o vermelho ordena em $K_0=2$)
    até a diagonal $t=\alpha$ perto de $(1,1)$.
    \item \textbf{M--F} (primeira ordem): em $T=0$, com $\sigma_i=\xi_i$ nos azuis e
    $\sigma_i=+1$ nos vermelhos, $m=q=\tfrac12$ e $f=-\tfrac18J_0-\tfrac14J$,
    contra $f=-\tfrac12J_0$ da fase F. Logo M--F em $\alpha=\tfrac23$.
    \item \textbf{SG--F} (primeira ordem, $t\gtrsim0{,}38$): obtida numericamente por
    comparação das energias livres das soluções de \eqref{eq:vh_sc4}; o salto
    de $m$ vai a zero em $(1,1)$.
\end{itemize}
%
O ponto $(1,1)$ é \textbf{triplo}: encontro de P, SG e F. Há ainda um
ponto triplo SG/M/F em $(\alpha,t)\approx(0{,}75;\,0{,}38)$.

\begin{figure}[H]
\centering
\begin{tikzpicture}[scale=4.2,>=stealth,font=\small]
  \draw[->] (0,0) -- (1.55,0) node[right] {$\alpha=J_0/J$};
  \draw[->] (0,0) -- (0,1.22) node[above] {$t=k_BT/J$};
  \foreach \x in {0.5,1,1.5}{\draw (\x,0.015)--(\x,-0.015) node[below]{$\x$};}
  \foreach \y in {0.5,1}{\draw (0.015,\y)--(-0.015,\y) node[left]{$\y$};}
  % P-SG e P-F
  \draw[very thick] (0,1) -- (1,1);
  \draw[very thick] (1,1) -- (1.2,1.2);
  % SG-M (contínua)
  \draw[thick,smooth] plot coordinates {(0,0) (0.080,0.040) (0.160,0.080) (0.240,0.120) (0.320,0.160) (0.400,0.200) (0.480,0.240) (0.558,0.280) (0.635,0.320) (0.691,0.350) (0.726,0.370) (0.749,0.383)};
  % M-F (1a ordem)
  \draw[thick,dashed,smooth] plot coordinates {(0.667,0) (0.667,0.050) (0.667,0.100) (0.669,0.150) (0.676,0.200) (0.690,0.250) (0.710,0.300) (0.728,0.340) (0.739,0.360) (0.749,0.383)};
  % SG-F (1a ordem)
  \draw[thick,dashed,smooth] plot coordinates {(0.749,0.383) (0.787,0.450) (0.836,0.550) (0.879,0.650) (0.918,0.750) (0.953,0.850) (0.976,0.920) (0.988,0.960) (1,1)};
  \fill (1,1) circle (0.4pt) node[above left] {$(1,1)$};
  \fill (0.749,0.383) circle (0.4pt);
  \node at (0.55,1.1) {P};
  \node at (0.35,0.62) {SG};
  \node at (0.50,0.10) {M};
  \node at (1.15,0.40) {F};
  \node at (1.30,1.30) [font=\scriptsize] {};
\end{tikzpicture}
\caption{Diagrama de fases do modelo de van Hemmen ($\xi,\eta=\pm1$), calculado
numericamente a partir de \eqref{eq:vh_sc4} e das energias livres. Linhas
cheias: transições contínuas; tracejadas: primeira ordem. $(1,1)$ é ponto triplo.}
\label{fig:vh_phase}
\end{figure}

\subsection*{Comparação com o modelo de Sherrington--Kirkpatrick}
O modelo SK tem o mesmo termo ferromagnético, mas acoplamentos aleatórios
independentes com $\overline{J_{ij}^2}=J^2/N$. Na aproximação RS,
%
\begin{equation}
    m=\int\!Dz\,\tanh\!\big[\beta(J_0m+J\sqrt q\,z)\big],\qquad
    q=\int\!Dz\,\tanh^2\!\big[\beta(J_0m+J\sqrt q\,z)\big].
\end{equation}
%
As instabilidades de P ($t=\alpha$ e $t=1$) coincidem com as do van Hemmen. A
fronteira RS SG--F ($m\to0$) é $\alpha=t\big/\!\int\!Dz\,\mathrm{sech}^2(\sqrt q\,z/t)$ e a linha de
Almeida--Thouless é $1=t^{-2}\!\int\!Dz\,\mathrm{sech}^4[(\alpha m+\sqrt q\,z)/t]$.

\begin{figure}[H]
\centering
\begin{tikzpicture}[scale=4.2,>=stealth,font=\small]
  \draw[->] (0,0) -- (1.55,0) node[right] {$\alpha=J_0/J$};
  \draw[->] (0,0) -- (0,1.22) node[above] {$t=k_BT/J$};
  \foreach \x in {0.5,1,1.5}{\draw (\x,0.015)--(\x,-0.015) node[below]{$\x$};}
  \foreach \y in {0.5,1}{\draw (0.015,\y)--(-0.015,\y) node[left]{$\y$};}
  \draw[very thick] (0,1) -- (1,1);
  \draw[very thick] (1,1) -- (1.2,1.2);
  % F-SG (RS, m->0)
  \draw[thick,smooth] plot coordinates {(1,1) (1.002,0.910) (1.008,0.830) (1.015,0.770) (1.026,0.690) (1.040,0.610) (1.053,0.550) (1.072,0.470) (1.089,0.410) (1.114,0.330) (1.143,0.250) (1.167,0.190) (1.201,0.110) (1.229,0.050)};
  % AT
  \draw[thick,dotted,smooth] plot coordinates {(1,1) (1.013,0.880) (1.029,0.820) (1.051,0.760) (1.079,0.700) (1.102,0.660) (1.140,0.600) (1.186,0.540) (1.239,0.480) (1.279,0.440) (1.348,0.380) (1.429,0.320) (1.526,0.260)};
  \fill (1,1) circle (0.4pt) node[above left] {$(1,1)$};
  \node at (0.5,0.5) {SG};
  \node at (1.40,0.15) {F+SG};
  \node at (1.40,0.75) {F};
\end{tikzpicture}
\caption{Diagrama do modelo SK (RS) no mesmo plano. Linha cheia: $m\to0$ (SG--F);
pontilhada: AT (a RS é instável abaixo dela). A SG é RSB para todo $\alpha<1$.}
\label{fig:sk_phase}
\end{figure}

\paragraph{Diferença.} No van Hemmen há só três parâmetros de ordem
$(m,q_1,q_2)$ e a solução é exata sem réplicas, com uma fase mista bem
localizada em baixo $\alpha$ e $T$ (o vermelho ordenando). No SK, a desordem
extensiva exige RSB: a região entre a linha RS SG--F e a AT é instável, e
a fronteira SG--F ocorre em $\alpha\gtrsim1$, não em $\alpha<1$ como no van Hemmen.
%
\begin{thebibliography}{99}
\bibitem{Aharony1978a}
  A.~Aharony,
  \textit{Tricritical points in systems with random fields},
  Phys.\ Rev.\ B \textbf{18}, 3318 (1978).

\bibitem{Aharony1978b}
  A.~Aharony,
  \textit{Spin-flop multicritical points in systems with random fields
  and in spin glasses},
  Phys.\ Rev.\ B \textbf{18}, 3328 (1978).

\bibitem{SchneiderPytte1977}
  T.~Schneider e E.~Pytte,
  \textit{Random-field instability of the ferromagnetic state},
  Phys.\ Rev.\ B \textbf{15}, 1519 (1977).

\bibitem{vanHemmen1982}
  J.~L.\ van Hemmen,
  \textit{Classical Spin-Glass Model},
  Phys.\ Rev.\ Lett.\ \textbf{49}, 409 (1982).
  
\end{thebibliography}

\end{document}
